\documentclass[12pt,a4paper]{article}
\usepackage[utf8]{inputenc}
\usepackage[T1]{fontenc}
\usepackage[left=1.4cm,right=1.4cm,top=0.95cm,bottom=0.95cm]{geometry}
\usepackage[hidelinks]{hyperref}
\usepackage{enumitem}
\usepackage{titlesec}
\pagestyle{empty}
\setlength{\parindent}{0pt}
\titleformat{\section}{\large\bfseries}{}{0em}{}[\titlerule]
\titlespacing*{\section}{0pt}{3.5pt}{2pt}
\setlist[itemize,1]{leftmargin=0pt,label={},topsep=0.5pt,itemsep=1.5pt,parsep=0pt}
\setlist[itemize,2]{leftmargin=1.2em,label=\textbullet,topsep=0pt,itemsep=0pt,parsep=0pt}

\begin{document}

{\centering
{\LARGE\textbf{Nicolas Rojas Varela}}\\[0.5mm]
{\large\textbf{Ingénieur R\&D en IA | Research Engineer}}\\[0.5mm]
Annecy, Auvergne-Rhône-Alpes, France\\
+33 7 45 61 10 20 \textbar{} \href{mailto:nicolasrojas1111@gmail.com}{nicolasrojas1111@gmail.com} \textbar{} \href{https://github.com/nirojasva}{GitHub}\\
Autorisé à travailler en France\par}
\vspace{1pt}

\section*{Profil}
Chercheur en informatique et ingénieur R\&D spécialisé dans l'apprentissage automatique appliqué aux séries temporelles, aux flux de données et à la détection d'anomalies. Expérience en conception d'algorithmes, évaluation expérimentale et recherche appliquée dans des contextes académique et industriel. Doctorant CIFRE à l'Université Clermont Auvergne, en collaboration avec Pfeiffer Vacuum+Fab Solutions, avec une soutenance prévue en mars 2027 (\href{https://spectral-anomaly-detector.limos.fr/index_en.html}{URL Page Web de Thèse}). Recherche un poste d'Ingénieur R\&D en IA ou de Research Engineer.

\section*{Expertise en recherche}
\textbf{Domaines :} Détection d'anomalies, apprentissage non supervisé, flux de données, séries temporelles, données spectrales, classification, interprétabilité\\
\textbf{Méthodes :} K-NN, Isolation Forest, autoencodeurs, évaluation comparative (benchmarking), reproductibilité expérimentale\\
\textbf{Programmation et bibliothèques :} Python, scikit-learn, CapyMOA, PySAD, Aeon, sktime, PyTorch, TensorFlow\\
\textbf{Outils complémentaires :} SQL, R, Git, GCP, pipelines de données, ETL, API

\section*{Expérience en R\&D et recherche}
\begin{itemize}
\item \textbf{Ingénieur R\&D / Doctorant CIFRE} \hfill Mars 2024 -- mars 2027 (prévu)\\
\textit{Pfeiffer Vacuum+Fab Solutions et LIMOS} \hfill Annecy, France
\begin{itemize}
\item Conception et évaluation de méthodes non supervisées pour la détection d'anomalies spectrales en temps réel dans un contexte industriel (modèle OnlineBootKNN : 85\% de taux de détection et 0,5\% de faux positifs).
\item Formulation d'approches d'apprentissage en ligne sur des flux de données spectrales, avec optimisation de l'efficacité algorithmique et de l'interprétabilité des modèles (temps de traitement inférieur à 1 s par spectre).
\item Développement de pipelines expérimentaux reproductibles en Python avec CapyMOA, scikit-learn et PySAD.
\end{itemize}

\item \textbf{Chercheur en apprentissage automatique} \hfill Novembre 2023 -- mars 2024\\
\textit{LIMOS} \hfill Aubière, France
\begin{itemize}
\item Développement d'algorithmes de classification de séries temporelles et de détection d'anomalies dans les flux de données (méthode interprétable RSAST à base de shapelets, temps d'exécution de 3 min 15 s sur des séries de longueur 1024).
\item Conception de protocoles expérimentaux comparatifs, évaluation de modèles et intégration open-source avec Aeon, scikit-learn et sktime.
\end{itemize}

\item \textbf{Stagiaire Data Miner} \hfill Mars 2023 -- septembre 2023\\
\textit{LIMOS} \hfill Aubière, France
\begin{itemize}
\item Analyse et classification interprétable de séries temporelles à l'aide de méthodes et de bibliothèques Python, notamment Aeon, scikit-learn et sktime.
\end{itemize}
\end{itemize}

\section*{Publications scientifiques}
\begin{itemize}
\item \textbf{Varela, N. R.}, Ah-Pine, J., \& Nguifo, E. M. (2026). \textit{OnlineBootKNN: An Unsupervised Framework for Detecting Anomalies in Spectral Data Streams}. \textit{Proceedings of the AAAI Conference on Artificial Intelligence}, 40(18), 15797--15805. \href{https://doi.org/10.1609/aaai.v40i18.38611}{DOI} \textbar{} \href{https://github.com/nirojasva/spectral-benchmark}{Code}
\item \textbf{Varela, N. R.}, Ah-Pine, J., \& Nguifo, E. M. (2026). \textit{The BV4 Benchmark for Unsupervised Anomaly Detection in High-Dimensional Spectral Data Streams}. \textit{Advances in Neural Information Processing Systems (NeurIPS), Datasets and Benchmarks Track}.
\item \textbf{Varela, N. R.}, Mbouopda, M. F., \& Nguifo, E. M. (2023). \textit{RSAST: Sampling Shapelets for Interpretable Time Series Classification}. \textit{Conférence Francophone sur l'Apprentissage Automatique (CAP)}. \href{https://hal.science/hal-04311309/}{HAL} \textbar{} \href{https://www.aeon-toolkit.org/en/v1.1.0/api_reference/auto_generated/aeon.transformations.collection.shapelet_based.RSAST.html}{Code (Aeon)}
\end{itemize}

\section*{Workshops et présentations}
\begin{itemize}
\item Amroun, C., \textbf{Varela, N. R.}, \& Nguifo, E. M. (2026). \textit{Enhanced OBKNN: An Efficient Framework for Anomaly Detection in Spectral Data Streams}. \textit{40th Conference on Neural Information Processing Systems (NeurIPS 2026), Workshop on Time Series in the Age of Large Models (TS-LIMITS)}. \href{https://hal.science/hal-05803665/}{HAL}
\item \textbf{Varela, N. R.}, Ah-Pine, J., \& Nguifo, E. M. (2026). \textit{OnlineBootKNN: An Unsupervised Framework for Detecting Anomalies in Spectral Data Streams}. \textit{ECML PKDD 2026 (Nectar Track)}.
\end{itemize}

\section*{Formation}
\begin{itemize}
\item \textbf{Doctorat en informatique} -- \textit{Université Clermont Auvergne} (Soutenance prévue en mars 2027) \hfill 2024 -- 2027
\item \textbf{Master international en informatique} -- \textit{ISIMA, Clermont-Ferrand} \hfill 2022 -- 2023
\item \textbf{Spécialisation en statistiques} -- \textit{Universidad Nacional de Colombia} \hfill 2020 -- 2021
\item \textbf{Ingénierie industrielle} -- \textit{Universidad Nacional de Colombia} \hfill 2012 -- 2017
\end{itemize}

\section*{Expérience professionnelle complémentaire}
\begin{itemize}
\item \textbf{Analyste Business Intelligence} \hfill Octobre 2021 -- août 2022\\
\textit{Center Source, Inc.} \hfill Barranquilla, Colombie
\begin{itemize}
\item Conception de pipelines ETL et modélisation de données (suivi d'indicateurs clients clés pour le business, actualisés en temps réel).
\item Développement de tableaux de bord et de solutions analytiques avec Python, R, Tableau et Exasol.
\end{itemize}

\item \textbf{Analyste de données} \hfill Août 2018 -- octobre 2021\\
\textit{Seguros Bolívar S.A.} \hfill Bogotá, Colombie
\begin{itemize}
\item Extraction, modélisation et analyse de données à grande échelle (amélioration du niveau de service client jusqu'à 80\%).
\item Utilisation de Python, R, Excel, GCP, Dataiku et Power BI en support aux activités opérationnelles.
\end{itemize}
\end{itemize}

\section*{Langues}
Espagnol : langue maternelle \textbar{} Anglais : niveau professionnel \textbar{} Français : niveau professionnel

\end{document}